% !TEX TS-program = pdfLaTeX+shellescape
% !TEX encoding = UTF-8 Unicode

\documentclass[crop=true]{standalone}
%\input{../../mypackages}
\usepackage{pgfplots}
\usepackage{multido}
\usepgfplotslibrary{fillbetween}

\begin{filecontents}{box.dat}
0.05  1.0515
0.15  1.1348
0.25  1.1667
0.35  1.1348
0.45  1.0515
0.55  0.9485
0.65  0.8652
0.75  0.8333
0.85  0.8652
0.95  0.9485
\end{filecontents}

\begin{document}
\begin{tikzpicture}
\begin{axis}[
axis lines = middle,
unit vector ratio*=1 1 1,
scale=1.2, xmin=-0.1, xmax=1.35, ymin=-0.2, ymax=0.35,
%grid = major,
xtick={0.2,1.2},
ytick={0},
xticklabels = {$a$, $b$},
yticklabels = {},
xlabel=$x$,
ylabel=$y$,
xlabel style={below right},
ylabel style={above left},
axis line style = {->},
]
\fill[gray!20] (axis cs:0.2,0) rectangle (axis cs:0.3,0.1515);
\draw (axis cs:0.2,0) rectangle (axis cs:0.3,0.1515);
\fill[gray!20] (axis cs:0.3,0) rectangle (axis cs:0.4,0.2348);
\draw (axis cs:0.3,0) rectangle (axis cs:0.4,0.2348);
\fill[gray!20] (axis cs:0.4,0) rectangle (axis cs:0.5,0.2667);
\draw (axis cs:0.4,0) rectangle (axis cs:0.5,0.2667);
\fill[gray!20] (axis cs:0.5,0) rectangle (axis cs:0.6,0.2348);
\draw (axis cs:0.5,0) rectangle (axis cs:0.6,0.2348);
\fill[gray!20] (axis cs:0.6,0) rectangle (axis cs:0.7,0.1515);
\draw (axis cs:0.6,0) rectangle (axis cs:0.7,0.1515);
\fill[gray!20] (axis cs:0.7,0) rectangle (axis cs:0.8,0.0485);
\draw (axis cs:0.7,0) rectangle (axis cs:0.8,0.0485);
\fill[gray!20] (axis cs:0.8,0) rectangle (axis cs:0.9,-0.0348);
\draw (axis cs:0.8,0) rectangle (axis cs:0.9,-0.0348);
\fill[gray!20] (axis cs:0.9,0) rectangle (axis cs:1.0,-0.0667);
\draw (axis cs:0.9,0) rectangle (axis cs:1.0,-0.0667);
\fill[gray!20] (axis cs:1.0,0) rectangle (axis cs:1.1,-0.0348);
\draw (axis cs:1.0,0) rectangle (axis cs:1.1,-0.0348);
\fill[gray!20] (axis cs:1.1,0) rectangle (axis cs:1.2,0.0485);
\draw (axis cs:1.1,0) rectangle (axis cs:1.2,0.0485);

\addplot[mark=none,thick] gnuplot[domain=0.2:1.2,samples=100]{(sin(2*pi*(x-0.2))+0.6)/6};

\end{axis}

\end{tikzpicture}


\end{document}